% ============================================================
% PHẦN SADGS — Slide 2/8: Structure Tensor & tỉ số tần số eta
% ============================================================

% --- Frame 1 ---
\begin{frame}{Đo cấu trúc texture đa tỉ lệ: Structure Tensor $\to$ $\eta$}
\footnotesize
\begin{itemize}
    \item Với mỗi camera nhìn thấy Gaussian, lấy mẫu \textbf{structure tensor} $S$ tại điểm chiếu $(u,v)$ trên ảnh gốc (được tiền tính đa tỉ lệ, \texttt{st\_levels} mức, \texttt{st\_mode}):
\end{itemize}
\[
S = \begin{pmatrix} S_{xx} & S_{xy} \\ S_{xy} & S_{yy}\end{pmatrix}, \qquad
\lambda_1 = \frac{\text{tr}(S)}{2} + \sqrt{\left(\frac{\text{tr}(S)}{2}\right)^2 - \det(S)}
\]
\begin{itemize}
    \item $\lambda_1$ là trị riêng lớn nhất (năng lượng gradient cực đại) $\Rightarrow$ \textbf{bước sóng cục bộ} $w_{\min} = 1/(\sqrt{\lambda_1}+\varepsilon)$ (đơn vị pixel) -- vùng chi tiết mịn có $w_{\min}$ nhỏ, vùng phẳng có $w_{\min}$ lớn.
    \item Chiều dài trục chiếu 2D của Gaussian (từ Jacobian phối cảnh $J$, ma trận xoay $R$, scale $\Sigma$): $\text{axis}_k = \lVert J R \,\text{diag}(\Sigma)\, e_k\rVert$, $k\in\{x,y,z\}$.
\end{itemize}
\[
\eta_k^{\text{wavelength}} = \frac{\text{axis}_k}{w_{\min}}, \qquad
\eta_k^{\text{projection}} = \sqrt{\,e_k^\top R^\top J^\top S J R\, e_k\,}
\]
\begin{itemize}
    \item Hai chế độ (\texttt{eta\_compute\_mode}): \texttt{"wavelength"} (mặc định, so kích thước Gaussian với bước sóng texture) hoặc \texttt{"projection"} (chiếu trực tiếp structure tensor lên từng trục).
    \item $\eta_k > 1$: trục $k$ của Gaussian \textbf{lớn hơn} chi tiết ảnh tại đó $\Rightarrow$ nguy cơ alias/mờ $\Rightarrow$ ứng viên split theo trục đó.
    \item $\eta$ được tính \textbf{riêng cho từng trục} ($\eta_{3ch}$, 3 kênh), không phải một số vô hướng chung cho cả Gaussian.
\end{itemize}
\end{frame}

% --- Frame 2 ---
\begin{frame}{Structure tensor đa tỉ lệ là gì?}
\begin{columns}
\begin{column}{0.55\textwidth}
\footnotesize
\begin{itemize}
    \item $S$ được xây từ \textbf{DoG (Difference of Gaussians)} đa octave trên ảnh ground-truth (\texttt{loss\_utils.py: get\_multiscale\_structure\_tensor\_v1}): structure tensor gốc $S_{\text{base}}$ (Di Zenzo + Sobel) được tích phân lại ở mỗi octave với $\rho_i = 3\sigma_i$ thành $S_i$, chuẩn hoá hướng $\hat S_i = S_i/(\mathrm{tr}\,S_i+\varepsilon)$ rồi cộng dồn có trọng số
    $S = \frac{1}{\sum_i w_i+\varepsilon}\sum_i \hat S_i\, w_i\, f_i^2$, với $w_i = \lVert I_{\sigma_i}-I_{\sigma_{i+1}}\rVert_2^{\,p}$ ($p=$\texttt{power\_factor}$=3$) và $f_i = 1/\texttt{octave\_step}^i$.
    \item Kết quả được \textbf{cache theo ảnh} và tra cứu lại trong \texttt{update\_freq\_stats\_online}, không tính lại mỗi iteration -- giữ chi phí runtime thấp.
    \item Trị riêng lớn nhất $\lambda_1$ của $S$ = hướng \& độ mạnh của biến thiên cường độ pixel mạnh nhất tại điểm đó $\Rightarrow$ chính là \textbf{tần số không gian cục bộ} cao nhất của texture.
    \item Vùng cạnh sắc/hoạ tiết dày đặc $\Rightarrow$ $\lambda_1$ lớn $\Rightarrow$ $w_{\min}$ nhỏ $\Rightarrow$ đòi hỏi Gaussian nhỏ để không làm mờ chi tiết.
    \item Vùng mảng màu phẳng (tường, bầu trời) $\Rightarrow$ $\lambda_1$ nhỏ $\Rightarrow$ $w_{\min}$ lớn $\Rightarrow$ Gaussian lớn vẫn đủ, không cần chi tiết thêm.
\end{itemize}
\end{column}
\begin{column}{0.45\textwidth}
\centering
\includegraphics[width=0.95\linewidth]{fig_07_threshold_sweep.png}
\captionof{figure}{\footnotesize Minh hoạ định tính: độ nhạy của phân loại tần số khi quét ngưỡng $\tau$}
\end{column}
\end{columns}
\end{frame}

% --- Frame 3 ---
\begin{frame}{Phân loại quan sát: high / mid / low}
\footnotesize
\begin{itemize}
    \item Mỗi \textbf{quan sát} (một Gaussian trong một view) được gán nhãn dựa trên $\eta_{\max}=\max_k \eta_k$ lấy trên 3 trục (\texttt{freq\_utils.py:401}), theo hai ngưỡng cứng:
\end{itemize}
\[
\text{label}(\eta_{\max}) =
\begin{cases}
\text{\textbf{high}} & \eta_{\max} > \tau_{\text{high}} = 1.0 \quad \text{(chi tiết vượt kích thước hiện tại} \Rightarrow \text{ứng viên split)} \\
\text{\textbf{low}} & \eta_{\max} \le \tau_{\text{low}} = 0.1 \quad \text{(Gaussian thừa chi tiết so với texture} \Rightarrow \text{ứng viên prune)} \\
\text{mid} & \tau_{\text{low}} < \eta_{\max} \le \tau_{\text{high}} \quad \text{(giữ nguyên)}
\end{cases}
\]
\vspace{0.15cm}
\centering
\begin{tabular}{lccccl}
\toprule
\textbf{Ví dụ} & $\eta_x$ & $\eta_y$ & $\eta_z$ & $\eta_{\max}$ & \textbf{Nhãn của quan sát} \\
\midrule
Gaussian A & 1.3 & 0.4 & 0.08 & 1.3 & \textbf{high} $\Rightarrow$ \texttt{eta\_high\_count} $+1$ \\
Gaussian B & 0.6 & 0.5 & 0.7 & 0.7 & mid $\Rightarrow$ \texttt{eta\_mid\_count} $+1$ \\
Gaussian C & 0.08 & 0.05 & 0.03 & 0.08 & \textbf{low} $\Rightarrow$ \texttt{eta\_low\_count} $+1$ \\
\bottomrule
\end{tabular}
\vspace{0.2cm}
\begin{itemize}
    \item Nhãn high/mid/low chỉ quyết định \textbf{có} split/prune hay không. \textbf{Mức độ} chia mới dùng $\eta$ \textbf{riêng từng trục}: \texttt{max\_eta\_3ch} (max qua các view, \texttt{freq\_utils.py:390--392}) $\to k_{\text{axis}}$ ở Slide~4.
    \item Ngưỡng lọc bổ sung \texttt{freq\_grad\_threshold} và \texttt{importance\_score\_threshold} ($=0.5$) giúp loại các tín hiệu tần số yếu/nhiễu trước khi đưa vào quyết định densify.
\end{itemize}
\end{frame}

% --- Frame 4 ---
\begin{frame}{Tính nhất quán đa góc nhìn (multiview-consistency)}
\begin{columns}
\begin{column}{0.55\textwidth}
\footnotesize
\begin{itemize}
    \item Một Gaussian được nhiều camera nhìn thấy trong quá trình huấn luyện; mỗi view cho một ước lượng $\eta_k$ khác nhau (do góc chiếu, occlusion, độ phân giải khác nhau).
    \item \texttt{update\_freq\_stats\_online} \textbf{tích luỹ} qua từng view vào bộ đếm trên mỗi Gaussian: \texttt{accum\_eta}, \texttt{eta\_high\_count}, \texttt{eta\_low\_count}, \texttt{accum\_view\_count}.
\end{itemize}
\[
\text{high\_ratio} = \frac{\texttt{eta\_high\_count}}{\texttt{accum\_view\_count}}, \qquad
\text{low\_ratio} = \frac{\texttt{eta\_low\_count}}{\texttt{accum\_view\_count}}
\]
\begin{itemize}
    \item Chỉ khi \text{high\_ratio} (hoặc \text{low\_ratio}) đủ lớn qua \textbf{nhiều view} thì Gaussian mới thực sự được coi là cần split (hoặc prune).
    \item Lý do: một view đơn lẻ có thể bị occlusion, mờ do chuyển động, hoặc góc nhìn xấu $\Rightarrow$ tín hiệu $\eta$ nhiễu tại view đó. Tích luỹ đa view giúp lọc nhiễu, tránh split/prune sai vì một khung hình bất thường.
\end{itemize}
\end{column}
\begin{column}{0.45\textwidth}
\centering
\includegraphics[width=0.95\linewidth]{fig_04_gbar_history.png}
\captionof{figure}{\footnotesize Minh hoạ khái niệm: giá trị tích luỹ ổn định dần qua nhiều lần cập nhật (hình gốc vẽ cho gradient trung bình, dùng ở đây để minh hoạ ý tưởng tích luỹ theo view, không phải đúng 100\% cơ chế $\eta$)}
\end{column}
\end{columns}
\end{frame}

% --- Frame 5 ---
\begin{frame}{So sánh hai chế độ \texttt{eta\_compute\_mode}}
\footnotesize
\centering
\begin{tabular}{lll}
\toprule
\textbf{Tiêu chí} & \textbf{"wavelength"} & \textbf{"projection"} \\
\midrule
Công thức & $\eta_k = \text{axis}_k / w_{\min}$ & $\eta_k = \sqrt{e_k^\top R^\top J^\top S J R\, e_k}$ \\
Ý nghĩa & So kích thước Gaussian & Chiếu trực tiếp dạng toàn phương \\
 & với 1 con số bước sóng vô hướng & của $S$ lên từng trục Gaussian \\
Thông tin dùng từ $S$ & Chỉ $\lambda_1$ (trị riêng lớn nhất) & Toàn bộ ma trận $S$ (cả hướng) \\
Độ nhạy hướng & Đẳng hướng (isotropic) theo $w_{\min}$ & Bất đẳng hướng, theo đúng hướng trục \\
Vai trò & Chế độ mặc định, đơn giản, ổn định & Lựa chọn thiết kế thay thế, tận dụng \\
 & & đầy đủ thông tin hướng của $S$ \\
\bottomrule
\end{tabular}
\vspace{0.3cm}
\begin{itemize}
    \item \textbf{"wavelength"}: quy $S$ về một đại lượng vô hướng ($\lambda_1 \to w_{\min}$) rồi so trực tiếp với độ dài trục -- đơn giản, dễ diễn giải như "kích thước Gaussian so với chi tiết ảnh".
    \item \textbf{"projection"}: giữ nguyên thông tin hướng của $S$, chiếu quadratic form lên từng trục $e_k$ -- có thể phân biệt tốt hơn khi texture có tính bất đẳng hướng mạnh (vd. vân gỗ, sọc kẻ).
    \item Đây là hai lựa chọn thiết kế đã cài sẵn trong \texttt{arguments/\_\_init\_\_.py} (\texttt{eta\_compute\_mode}); code không kèm ghi chú lý do chọn mặc định, nên không suy diễn thêm ngoài khác biệt công thức nêu trên.
\end{itemize}
\end{frame}
